\documentclass[12pt,a4paper]{article}
\usepackage[utf8]{inputenc}
\usepackage{amsmath}
\usepackage{graphicx}
\usepackage{geometry}
\usepackage{booktabs}
\geometry{a4paper, margin=1in}

\title{Vehicle Routing Optimization Report: Mathematical Comparison of Fleet Size}
\author{Logistics AI Agent}
\date{\today}

\begin{document}

\maketitle

\section*{1. Executive Summary}
This report provides a strict mathematical and operational comparison between different fleet size strategies for delivering a generated dataset of 50 packages over the Greater Tunis region. We analyze the scenarios of forcing 1 single vehicle, using the AI's optimized suggestion of 2 vehicles, and dividing the labor evenly across 5 vehicles. 

\section*{2. The Mathematical Model}
To determine the ``best'' number of vehicles, we utilize a unified cost function ($Z$) comprising the variable cost of route traversal (fuel, wear and tear) and the fixed cost of deploying a vehicle (driver base salary, rental/amortization).

$$ Z = \sum_{v=1}^{V} \left( C_{fixe} + C_{km} \times d_v \right) $$

Where:
\begin{itemize}
    \item $V$ = Number of vehicles dispatched.
    \item $C_{fixe}$ = Fixed cost per vehicle per day (Estimated at 150 TND for salary + dispatch overhead).
    \item $C_{km}$ = Variable cost per kilometer (Estimated at 0.20 TND/km for fuel and maintenance).
    \item $d_v$ = Total distance traveled by vehicle $v$.
\end{itemize}

Additionally, we impose a hard constraint such that $d_v \le 200 \text{ km}$ to ensure a driver can physically complete the route within an 8-hour workday, considering Tunis traffic and delivery handling times.

\section*{3. Scenario Analysis}

\subsection*{Scenario 1: Using Exactly 1 Vehicle}
\begin{itemize}
    \item \textbf{Total Distance traversed:} 252.61 km
    \item \textbf{Calculation:} $Z_1 = (1 \times 150) + (252.61 \times 0.20) = 200.52 \text{ TND}$
    \item \textbf{Constraint Check:} \textbf{FAILED.} The distance $d_1 = 252.61 \text{ km}$ strictly violates the logistical upper bounds of a single workday. Dealing with 50 packages spread across 120km territory is impossible for one entity.
\end{itemize}

\subsection*{Scenario 2: The Optimized Choice (2 Vehicles)}
By giving the OR-Tools algorithm the freedom to choose while penalizing unnecessary vehicle spawns, the solver naturally converges on \textbf{2 vehicles}.
\begin{itemize}
    \item \textbf{Total Distance traversed:} 303.45 km
    \item \textbf{Distribution:} Vehicle 1 (198.08 km) \& Vehicle 2 (105.37 km).
    \item \textbf{Calculation:} $Z_2 = (2 \times 150) + (303.45 \times 0.20) = 300 + 60.69 = 360.69 \text{ TND}$
    \item \textbf{Constraint Check:} \textbf{PASSED.} All routes fall within acceptable human and logistical limits.
\end{itemize}

\subsection*{Scenario 3: Using 5 Vehicles (Even spread)}
If we ignore fixed costs and merely try to balance the work among 5 drivers:
\begin{itemize}
    \item \textbf{Total Distance traversed:} 371.13 km (Notice the spike in km due to multiple round-trips to the central depot).
    \item \textbf{Distribution:} Routes average $\approx 74 \text{ km}$ per driver.
    \item \textbf{Calculation:} $Z_5 = (5 \times 150) + (371.13 \times 0.20) = 750 + 74.22 = 824.22 \text{ TND}$
    \item \textbf{Constraint Check:} Passed, but terribly inefficient.
\end{itemize}

\section*{4. Conclusion and Metric Comparisons}

\begin{table}[h]
\centering
\begin{tabular}{@{}lcccc@{}}
\toprule
\textbf{Scenario} & \textbf{Vehicles Used} & \textbf{Total Distance} & \textbf{Total Cost (TND)} & \textbf{Feasibility} \\ \midrule
Single Vehicle & 1 & 252.61 km & 200.52 & Impossible \\
\textbf{Optimized Choice} & \textbf{2} & \textbf{303.45 km} & \textbf{360.69} & \textbf{Optimal} \\
Max Spread Fleet & 5 & 371.13 km & 824.22 & Massive Loss \\ \bottomrule
\end{tabular}
\caption{Comparison of Logistical Optimization Metrics}
\end{table}

\textbf{Conclusion:} 
While sending only 1 vehicle technically results in the shortest cumulative GPS distance (by eliminating redundant depot return trips), it exceeds the physical capacity limit of what a human can deliver in one day. On the opposite spectrum, sending 5 vehicles results in wasted routing (a 22\% increase in total distance driven) and exorbitant fixed costs. The AI correctly identified that balancing exactly \textbf{2 vehicles} allows for functional execution while maintaining maximum possible profit margins.

\end{document}